\chapter{Conclusion}\label{chapter:conclusion}
Short introduction. It should relect on the primary and secondary research questions.

\section{Highlights}\label{sec:con-highlights}
Mention the strengths (such as adaptiveness through online learning) and weaknesses (such as more CPU time required) of the framework.

Also elaborate on the main contributions of the thesis.

\section{Limitations}\label{sec:stu-limitations}
Lay out the limitations of the study. It should include the difficulty of comparing single-threaded monolithic solvers to the parallelized framework (CPU time vs. wallclock time, different hardware, etc.).

Also mention limitations and variations during the run when it comes to harware and number of seeds. Computation was a limiting factor here and definitely influenced performance.

\section{Future Research}\label{sec:con-future}
Use the HAOS weights to train an offline learning model and inject the results back into the framework. This effectively eliminates the warm-up period and will lead to a substential speed-up of the framework.

Also mention that the framework, other than monolithic solvers, is highlz adaptable to different instances and VRP variants. Since it only requires the plug-and-play of capable clustering methods andsolvers on the their respective stages.